\section{Exercício Prático: Servidor Nginx}

\begin{frame}
  \frametitle{Exercício: Nginx}
  \textbf{Objetivo}: Subir um servidor Nginx, alterar a página inicial via terminal
  e monitorar os logs em tempo real.
\end{frame}

\begin{frame}
  \frametitle{1. Subindo o Servidor}
  \texttt{docker run -d -p 8080:80 --name meu-nginx nginx}

  \vspace{0.8em}
  \begin{itemize}
    \item Inicia o Nginx em segundo plano
    \item Porta 8080 do host $\rightarrow$ porta 80 do container
    \item Acessar \texttt{http://localhost:8080} exibe a página padrão
  \end{itemize}
\end{frame}

\begin{frame}[fragile]
  \frametitle{2. Alterando a Página via Terminal}
  Usando \texttt{docker exec} para acessar o shell do container:

  \vspace{0.5em}
\begin{verbatim}
docker exec -it meu-nginx bash

# Dentro do container:
echo '<h1>Minha pagina Docker!</h1>' \
  > /usr/share/nginx/html/index.html
exit
\end{verbatim}

  \vspace{0.5em}
  Recarregar \texttt{http://localhost:8080} mostra a nova página.

  \vspace{0.3em}
  \small O diretório \texttt{/usr/share/nginx/html/} é onde o Nginx serve arquivos estáticos.
\end{frame}

\begin{frame}
  \frametitle{3. Monitorando Logs em Tempo Real}
  \texttt{docker logs -f meu-nginx}

  \vspace{0.5em}
  \begin{itemize}
    \item \texttt{-f} segue os logs em tempo real
    \item Cada acesso no navegador gera uma nova linha de log
    \item \texttt{Ctrl+C} para sair
  \end{itemize}

  \vspace{0.5em}
  Ver apenas as últimas linhas:

  \vspace{0.3em}
  \texttt{docker logs --tail 10 meu-nginx}

  \vspace{0.8em}
  \textbf{Limpeza}: \texttt{docker rm -f meu-nginx}
\end{frame}
